\documentclass[11pt]{article}
\usepackage[a4paper,margin=27mm]{geometry}
\usepackage{amsmath,amssymb,amsthm,hyperref}
\newtheorem{theorem}{Theorem}
\newtheorem{lemma}[theorem]{Lemma}
\newtheorem{corollary}[theorem]{Corollary}
\title{Nearly linear face counts for exact three- and four-wise fairness}
\author{A consequence of permutation designs and palindrome composition}
\date{30 September 2026}
\begin{document}
\maketitle
Let $M_k(n)$ be the least common number of faces of $n$ uniform dice,
with all labels distinct, for which every subset of at most $k$ dice
has uniform relative order. Here $k$ is fixed while $n$ grows.
\begin{theorem}\label{main}
For $n\ge5$,
\[
 \frac{n+3}{2}\le M_3(n)\le M_4(n)\le Cn(\log n)^2
\]
for an absolute constant $C$. More precisely,
\[
 M_3(n)\le (6\sqrt e+o(1))n(\log_2n)^2,\qquad
 M_4(n)\le (12\sqrt e+o(1))n(\log_2n)^2.
\]
Thus $M_3(n)=n^{1+o(1)}$ and $M_4(n)=n^{1+o(1)}$.
\end{theorem}
The upper bounds use published permutation-design constructions; they
are not claims of new bounds for those designs. The lower bound is the
independent Gram/path-energy argument in
\path{../approximate/three_wise_face_lower_bound.tex}.

\begin{lemma}[A triple-minimum design gives three-wise fair dice]\label{pal}
Let $\rho_1,\ldots,\rho_q$ be permutations written as ordered lists of
$[n]$. Suppose that, for every three distinct labels, each is the earliest
of the three in exactly $q/3$ of these lists. Then
\[
 W=P_1\cdots P_q,\qquad P_i=\rho_i\operatorname{rev}(\rho_i),
\]
is three-wise fair and has $2q$ copies of each letter. The order of its
$q$ blocks is arbitrary.
\end{lemma}
\begin{proof}
Each palindrome block has two copies of each label and exactly two
occurrences of either ordering of any pair. For a target triple $abc$,
its count in $P_i$ is zero if $b$ is earliest in $\rho_i$, and two
otherwise. Thus the sum of its within-block counts is $4q/3$,
independent of $abc$.
The contributions using two distinct blocks depend only on their
one- and two-letter counts, and those using three distinct blocks depend
only on their one-letter counts. They too are independent of the target
order. Hence all six orders of any triple have the same count.
The one- and two-letter statements follow in the same way.
\end{proof}

\begin{lemma}[Reversal raises odd strength by one]\label{reverse}
If a word $W$ is fair on every subset of at most $2s-1$ labels, then
$W\operatorname{rev}(W)$ is fair on every subset of at most $2s$ labels.
Equal composition is not needed. Every multiplicity doubles.
\end{lemma}
\begin{proof}
Work in the associative algebra in which words with repeated letters
vanish, truncated above degree $2s$. Put $H=\sum_i c_iX_i$, where $c_i$
are the letter counts. Its positive-word signature is the product of
$1+X_i=\exp X_i$ in word order. The hypothesis says
\[
 \log S(W)=H+E_{2s}
\]
in this truncated algebra. The logarithm is a Lie element, since it is
obtained by the Baker--Campbell--Hausdorff formula.
Reversal acts on a homogeneous Lie element of degree $j$ by
$(-1)^{j-1}$, as follows by induction on brackets. Thus
\[
 \log S(\operatorname{rev}(W))=H-E_{2s}.
\]
All commutators involving $E_{2s}$ have degree greater than $2s$.
Consequently $\log S(W\operatorname{rev}(W))=2H$, precisely the
signature condition for fairness through degree $2s$.
\end{proof}

\begin{proof}[Proof of Theorem~\ref{main}]
A perfect sequence covering array of strength three is a multiset of
$q$ permutations in which every ordering of any triple occurs $q/6$
times. In particular it satisfies Lemma~\ref{pal}. Published
constructions give
\[
 q\le(3\sqrt e+o(1))n(\log_2n)^2;
\]
see Iurlano \cite[Remark 5, equation (11)]{iurlano}, using the finite-
geometry construction of Tarui, Itoh, and Takei with the improved base
case. Lemma~\ref{pal} gives $2q$ faces for three-wise fairness, and
Lemma~\ref{reverse} gives $4q$ for four-wise fairness.
The three-wise Gram/path-energy lower bound cited above gives
$M_3(n)\ge(n+3)/2$ for $n\ge5$.
\end{proof}

\paragraph{Questions exposed by the reduction.}
Can the logarithmic factor be removed for three-wise fair dice?
Does $M_3(n)=\Theta(n)$ hold, perhaps even without a linear-size
permutation design? The same question for $M_4(n)$ is equivalent up to
a factor of two by Lemma~\ref{reverse}. More generally that lemma gives
$M_{2s-1}(n)\le M_{2s}(n)\le2M_{2s-1}(n)$.
The next distinct regime is five-wise fairness. It is natural to ask
whether $M_5(n)$ is nearly quadratic, but no such lower or upper bound
is asserted here. These fixed-strength questions do not settle the
full-strength case $k=n$ in the manuscript.

\begin{thebibliography}{9}
\bibitem{iurlano} E.~Iurlano, \emph{Growth of the perfect sequence
covering array number}, Designs, Codes and Cryptography 91 (2023),
1487--1494,
\href{https://doi.org/10.1007/s10623-022-01168-3}{doi:10.1007/s10623-022-01168-3}.
\bibitem{yuster} R.~Yuster, \emph{Perfect sequence covering arrays},
Designs, Codes and Cryptography 88 (2020), 585--593;
\href{https://math.haifa.ac.il/raphy/papers/perfect-sca.pdf}{author's text}.
\end{thebibliography}
\end{document}
